% GENERATED FILE. Edit canonical lessons, then run npm run generate:books.
% canonical-source-hash: fnv1a64:6192a20608be5b14
% canonical-lessons: BN-C443-khobh, BN-C443-khitkhite, BN-C443-khepano, BN-C443-phuse-otha, BN-C443-gojgoj-kora

\chapter{Temper and Grumbling}
\label{ch:bn-temper-and-grumbling}
\hlchaptermodality{443}

\begin{chapteropening}
\textbf{By the end of this chapter:} \emph{I can say that someone is irritable or has lost their temper.}

Complete the last lesson of chapter 443: I can say that someone is irritable or has lost their temper.
\end{chapteropening}

\section[\texorpdfstring{\bn{ক্ষোভ} (khobh)}{ক্ষোভ (khobh)}]{\bn{ক্ষোভ} (khobh) — resentment; bitterness}

\label{lesson:BN-C443-khobh}



\begin{quote}
Before the new one: say the Bengali for free of worry, then the Bengali for to fidget.
\end{quote}

\subsection*{You'll want to know: \bn{ক্ষোভ}}
\textbf{\bn{ক্ষোভ}} — \emph{khobh} — \textquotedblleft{}resentment; bitterness\textquotedblright{}.

\begin{quote}
\bn{নতুন} \bn{নিয়ম} \bn{নিয়ে} \bn{কর্মীরা} \bn{ক্ষোভ} \bn{জানাল।}
\end{quote}

\begin{quote}
\emph{nôtun niyôm niye kôrmīrā khobh jānālô.}
\end{quote}

\begin{quote}
The workers voiced their resentment over the new rule.
\end{quote}

\begin{cousinweb}
\textbf{\bn{ক্ষোভ}} (\emph{khobh}) is a tatsama, Sanskrit \emph{kṣobha}, \textquotedblleft{}shaking, agitation, emotion\textquotedblright{} (Monier-Williams); in Bengali it is the stirred-up bitterness of someone who feels wronged.
\end{cousinweb}

\begin{grammarlens}[title={What you've built}]
One more word for describing what you lived and felt.
\end{grammarlens}

\subsection*{Your turn}
\begin{itemize}
\raggedright
  \item \emph{Say it:} \emph{khobh}
  \item \emph{Say it:} \emph{khobh}, once more
  \item \emph{Say it:} \emph{nôtun niyôm niye kôrmīrā khobh jānālô.} — The workers voiced their resentment over the new rule.
  \item \emph{Recall:} say the Bengali for free of worry, then the Bengali for to fidget, then say \emph{khobh} again
\end{itemize}

\begin{tcolorbox}[breakable,colback=teal!4,colframe=teal!35!black,title={Before you move on}]
What does \bn{ক্ষোভ} mean? (\textquotedblleft{}Resentment; bitterness\textquotedblright{}.) Say it once more.
\end{tcolorbox}
\section[\texorpdfstring{\bn{খিটখিটে} (khiṭkhiṭe)}{খিটখিটে (khiṭkhiṭe)}]{\bn{খিটখিটে} (khiṭkhiṭe) — irritable; crotchety}

\label{lesson:BN-C443-khitkhite}



\begin{quote}
Before the new one: say the Bengali for to fidget, then the Bengali for resentment.
\end{quote}

\subsection*{You'll want to know: \bn{খিটখিটে}}
\textbf{\bn{খিটখিটে}} — \emph{khiṭkhiṭe} — \textquotedblleft{}irritable; crotchety\textquotedblright{}.

\begin{quote}
\bn{জ্বর} \bn{হলে} \bn{দাদু} \bn{খিটখিটে} \bn{হয়ে} \bn{যান।}
\end{quote}

\begin{quote}
\emph{jbôr hôle dādu khiṭkhiṭe hôye jān.}
\end{quote}

\begin{quote}
When he has a fever, my grandfather gets crotchety.
\end{quote}

\begin{cousinweb}
\textbf{\bn{খিটখিটে}} (\emph{khiṭkhiṭe}) is an expressive word, its sound standing for a scratchy, easily annoyed temper; its origin is not settled.
\end{cousinweb}

\begin{grammarlens}[title={What you've built}]
One more word for describing what you lived and felt.
\end{grammarlens}

\subsection*{Your turn}
\begin{itemize}
\raggedright
  \item \emph{Say it:} \emph{khiṭkhiṭe}
  \item \emph{Say it:} \emph{khiṭkhiṭe}, once more
  \item \emph{Say it:} \emph{jbôr hôle dādu khiṭkhiṭe hôye jān.} — When he has a fever, my grandfather gets crotchety.
  \item \emph{Recall:} say the Bengali for to fidget, then the Bengali for resentment, then say \emph{khiṭkhiṭe} again
\end{itemize}

\begin{tcolorbox}[breakable,colback=teal!4,colframe=teal!35!black,title={Before you move on}]
What does \bn{খিটখিটে} mean? (\textquotedblleft{}Irritable; crotchety\textquotedblright{}.) Say it once more.
\end{tcolorbox}
\section[\texorpdfstring{\bn{খেপানো} (khepāno)}{খেপানো (khepāno)}]{\bn{খেপানো} (khepāno) — to provoke; to tease until someone gets cross}

\label{lesson:BN-C443-khepano}



\begin{quote}
Before the new one: say the Bengali for resentment, then the Bengali for irritable.
\end{quote}

\subsection*{You'll want to know: \bn{খেপানো}}
\textbf{\bn{খেপানো}} — \emph{khepāno} — \textquotedblleft{}to provoke; to tease until someone gets cross\textquotedblright{}.

\begin{quote}
\bn{পাড়ার} \bn{ছেলেরা} \bn{কুকুরটাকে} \bn{খেপাল।}
\end{quote}

\begin{quote}
\emph{pāṛār chhelerā kukurṭāke khepālô.}
\end{quote}

\begin{quote}
The neighbourhood boys provoked the dog.
\end{quote}

Past (he/she …): \bn{তিনি} \bn{ভাইকে} \bn{খেপালেন}, \emph{tini bhāike khepālen}.

Its plain partner is \bn{খেপে} \bn{যাওয়া}, to fly into a rage.

\begin{cousinweb}
\textbf{\bn{খেপানো}} (\emph{khepāno}) is the causative of \emph{khepā}, \textquotedblleft{}to go mad\textquotedblright{}, which is probably Sanskrit \emph{kṣipta}, \textquotedblleft{}thrown\textquotedblright{} (CDIAL 3685), a word Monier-Williams also gives for a scattered, distracted mind: to drive someone wild.
\end{cousinweb}

\begin{grammarlens}[title={What you've built}]
One more word for describing what you lived and felt.
\end{grammarlens}

\subsection*{Your turn}
\begin{itemize}
\raggedright
  \item \emph{Say it:} \emph{khepāno}
  \item \emph{Say it:} \emph{khepāno}, once more
  \item \emph{Say it:} \emph{pāṛār chhelerā kukurṭāke khepālô.} — The neighbourhood boys provoked the dog.
  \item \emph{Recall:} say the Bengali for resentment, then the Bengali for irritable, then say \emph{khepāno} again
\end{itemize}

\begin{tcolorbox}[breakable,colback=teal!4,colframe=teal!35!black,title={Before you move on}]
What does \bn{খেপানো} mean? (\textquotedblleft{}To provoke; to tease until someone gets cross\textquotedblright{}.) Say it once more.
\end{tcolorbox}
\section[\texorpdfstring{\bn{ফুঁসে} \bn{ওঠা} (phũse oṭhā)}{ফুঁসে ওঠা (phũse oṭhā)}]{\bn{ফুঁসে} \bn{ওঠা} (phũse oṭhā) — to seethe; to bristle with anger}

\label{lesson:BN-C443-phuse-otha}



\begin{quote}
Before the new one: say the Bengali for irritable, then the Bengali for to provoke.
\end{quote}

\subsection*{You'll want to know: \bn{ফুঁসে} \bn{ওঠা}}
\textbf{\bn{ফুঁসে} \bn{ওঠা}} — \emph{phũse oṭhā} — \textquotedblleft{}to seethe; to bristle with anger\textquotedblright{}.

\begin{quote}
\bn{খবরটা} \bn{শুনে} \bn{গ্রামের} \bn{লোকজন} \bn{ফুঁসে} \bn{উঠল।}
\end{quote}

\begin{quote}
\emph{khôbôrṭā shune grāmer lokjôn phũse uṭhlô.}
\end{quote}

\begin{quote}
When they heard the news, the villagers seethed.
\end{quote}

Past (he/she …): \bn{তিনি} \bn{রাগে} \bn{ফুঁসে} \bn{উঠলেন}, \emph{tini rāge phũse uṭhlen}.

\begin{cousinweb}
\textbf{\bn{ফুঁসে} \bn{ওঠা}} (\emph{phũse oṭhā}) is \emph{phũse}, the -\bn{এ} form of an imitative verb for the angry hiss of a snake, with \emph{oṭhā}, \textquotedblleft{}to rise\textquotedblright{} (CDIAL 1793).
\end{cousinweb}

\begin{grammarlens}[title={What you've built}]
One more word for describing what you lived and felt.
\end{grammarlens}

\subsection*{Your turn}
\begin{itemize}
\raggedright
  \item \emph{Say it:} \emph{phũse oṭhā}
  \item \emph{Say it:} \emph{phũse oṭhā}, once more
  \item \emph{Say it:} \emph{khôbôrṭā shune grāmer lokjôn phũse uṭhlô.} — When they heard the news, the villagers seethed.
  \item \emph{Recall:} say the Bengali for irritable, then the Bengali for to provoke, then say \emph{phũse oṭhā} again
\end{itemize}

\begin{tcolorbox}[breakable,colback=teal!4,colframe=teal!35!black,title={Before you move on}]
What does \bn{ফুঁসে} \bn{ওঠা} mean? (\textquotedblleft{}To seethe; to bristle with anger\textquotedblright{}.) Say it once more.
\end{tcolorbox}
\section[\texorpdfstring{\bn{গজগজ} \bn{করা} (gôjgôj kôrā)}{গজগজ করা (gôjgôj kôrā)}]{\bn{গজগজ} \bn{করা} (gôjgôj kôrā) — to grumble; to mutter crossly}

\label{lesson:BN-C443-gojgoj-kora}



\begin{quote}
Before the new one: say the Bengali for to provoke, then the Bengali for to seethe.
\end{quote}

\subsection*{You'll want to know: \bn{গজগজ} \bn{করা}}
\textbf{\bn{গজগজ} \bn{করা}} — \emph{gôjgôj kôrā} — \textquotedblleft{}to grumble; to mutter crossly\textquotedblright{}.

\begin{quote}
\bn{দাদু} \bn{সকাল} \bn{থেকে} \bn{গজগজ} \bn{করছেন।}
\end{quote}

\begin{quote}
\emph{dādu sôkāl theke gôjgôj kôrchhen.}
\end{quote}

\begin{quote}
Grandfather has been grumbling since the morning.
\end{quote}

Past (he/she …): \bn{তিনি} \bn{সারাদিন} \bn{গজগজ} \bn{করলেন}, \emph{tini sārādin gôjgôj kôrlen}.

\begin{cousinweb}
\textbf{\bn{গজগজ} \bn{করা}} (\emph{gôjgôj kôrā}) is \emph{gôjgôj}, an imitative word for low, cross muttering, with \emph{kôrā}, \textquotedblleft{}to do\textquotedblright{} (CDIAL 2814).
\end{cousinweb}

\begin{grammarlens}[title={What you've built}]
One more word for describing what you lived and felt.
\end{grammarlens}

\subsection*{Your turn}
\begin{itemize}
\raggedright
  \item \emph{Say it:} \emph{gôjgôj kôrā}
  \item \emph{Say it:} \emph{gôjgôj kôrā}, once more
  \item \emph{Say it:} \emph{dādu sôkāl theke gôjgôj kôrchhen.} — Grandfather has been grumbling since the morning.
  \item \emph{Recall:} say the Bengali for to provoke, then the Bengali for to seethe, then say \emph{gôjgôj kôrā} again
\end{itemize}

\begin{tcolorbox}[breakable,colback=teal!4,colframe=teal!35!black,title={Before you move on}]
What does \bn{গজগজ} \bn{করা} mean? (\textquotedblleft{}To grumble; to mutter crossly\textquotedblright{}.) Say it once more.
\end{tcolorbox}
